\section{Model and estimation targets}
\label{sec:model}

\subsection{Observed data}
\label{subsec:observations}
Bold letters denote vectors and matrices; indexed scalar entries are italic. For independent patients, we observe
\begin{equation}
 \begin{gathered}
O_i=(Y_i,\Delta_i,\bm{X}_i,\bm{Z}_i),\\
 Y_i=\log\{\min(T_i,C_i)\},\\
 \Delta_i=\mathbf 1\{T_i\le C_i\},\\
 i=1,\ldots,n,
 \end{gathered}
 \label{eq:observations}
\end{equation}
where $T_i\in(0,\infty)$ is the event time, $C_i\in(0,\infty]$ is the censoring time, $\bm{X}_i\in\mathbb R^p$ contains baseline molecular measurements, and $\bm{Z}_i\in\mathbb R^d$ contains prespecified clinical covariates. We consider a single right-censored endpoint observed from time zero, with fully observed baseline covariates and $Y_i$ on the natural-logarithmic scale.

The covariates have prespecified coordinate scales and finite second moments. We absorb population means into the unspecified error location and take $E(\bm{Z}_i)=E(\bm{X}_i)=0$ for notation. The clinical design excludes a constant column; Section~\ref{subsec:centering} defines sample centering. The dimension $p=p_n$ and sparsity may grow with $n$, whereas $d$ and the true factor dimension $K$ are fixed in the theory. Dependence of the joint law on $n$ is suppressed.

\subsection{Factor representation and coordinate normalization}
\label{subsec:factors}
Suppose the molecular covariates admit the representation
\begin{equation}
 \begin{gathered}
\bm{X}_i=\bm{B}\bm{f}_i+\bm{u}_i,\\
 \bm{B}\in\mathbb R^{p\times K},\quad
 \bm{f}_i\in\mathbb R^K,\quad \bm{u}_i\in\mathbb R^p,
 \end{gathered}
 \label{eq:factor-model}
\end{equation}
with
\begin{equation}
 E(\bm{f}_i)=0,\quad E(\bm{u}_i)=0,\quad
 E(\bm{f}_i\bm{f}_i^\top)=\bm{I}_K,\quad E(\bm{f}_i\bm{u}_i^\top)=0.
 \label{eq:factor-normalization}
\end{equation}
The factors and idiosyncratic components are latent. Define the full design vector
\[
 \bm{D}_i=(\bm{Z}_i^\top,\bm{f}_i^\top,\bm{u}_i^\top)^\top\in\mathbb R^P,
 \qquad q=d+K,\qquad P=p+q.
\]
Equation~\eqref{eq:factor-normalization} imposes second-moment orthogonality between factors and idiosyncratic components, while allowing dependence within and between covariate blocks. Clinical covariates may be associated with either component.

For the pervasive-factor regime, the basic structural conditions are
\begin{equation}
 \begin{gathered}
c_B \bm{I}_K\preceq p^{-1}\bm{B}^\top \bm{B}\preceq C_B \bm{I}_K,
 \\
 \|\bm{\Sigma}_u\|_{\mathrm{op}}\le C_u,
 \\
 \bm{\Sigma}_u=\operatorname{Cov}(\bm{u}_i),
 \end{gathered}
 \label{eq:pervasive}
\end{equation}
where the positive constants are uniform in $n,p$. Equation~\eqref{eq:pervasive} supplies the spectral separation used in factor-adjusted regression \citep{FanEtAl2020}; Section~\ref{sec:theory} adds distributional, loading, and sparse-design conditions for estimated-factor guarantees.

For an orthogonal $K\times K$ matrix $\bm{Q}$, the transformation
\begin{equation}
 \bm{f}_i^*=\bm{Q}\bm{f}_i,\qquad \bm{B}^*=\bm{B}\bm{Q}^\top
 \label{eq:factor-rotation}
\end{equation}
preserves the common component and~\eqref{eq:factor-normalization}; the factor coefficients transform with the same rotation. Population rotation and sample normalization $\widehat{\bm{F}}^\top\widehat{\bm{F}}/n=\bm{I}_K$ play separate roles. Finite-sample alignment matrices can be nonorthogonal \citep{BaiNg2013}. Section~\ref{subsec:identification} establishes slope identification for these coordinate representations.

\subsection{AFT models and the sparse coefficient target}
\label{subsec:aft-models}
The primary model is an ordinary sparse AFT model,
\begin{equation}
 \mathcal M_0:\qquad
 \log T_i=\bm{Z}_i^\top\bm{\theta}_0+\bm{X}_i^\top\bm{\beta}_0+\varepsilon_i.
 \label{eq:ordinary-aft}
\end{equation}
Factor adjustment retains $\bm{\beta}_0$ as the observed-covariate slope and represents dependence through common and residual components. We also consider the extension
\begin{equation}
\begin{aligned}
\mathcal M_1:\qquad\log T_i&=\eta_i^0+\varepsilon_i,\\
\eta_i^0&=\bm Z_i^\top\bm\theta_0+\bm X_i^\top\bm\beta_0+\bm f_i^\top\bm\xi_0\\
&=\bm Z_i^\top\bm\theta_0+\bm u_i^\top\bm\beta_0+\bm f_i^\top\bm\gamma_0,\\
\bm\gamma_0&=\bm B^\top\bm\beta_0+\bm\xi_0.
\end{aligned}
\label{eq:augmented-aft}
\end{equation}
Thus $\mathcal M_0$ is the submodel $\bm{\xi}_0=0$. Additional factor contributions also arise in factor-augmented hazard regression (Bayle and Fan, 2022, preprint, arXiv:2210.01067). In $\mathcal M_1$, the coefficient target is defined relative to the specified latent factor representation.

Write
\begin{equation}
 \begin{gathered}
\bm{a}_0=(\bm{\theta}_0^\top,\bm{\gamma}_0^\top)^\top,\\
 \bm{\vartheta}_0=(\bm{a}_0^\top,\bm{\beta}_0^\top)^\top,\\
 S=\operatorname{supp}(\bm{\beta}_0),\\
 s=|S|,\\
 k=s+q.
 \end{gathered}
 \label{eq:parameter-notation}
\end{equation}
Sparsity is imposed on $\bm{\beta}_0$; the $q$ clinical and factor coefficients are unpenalized. The target $\bm{\beta}_0$, its support, and $\eta_i^0$ are invariant under~\eqref{eq:factor-rotation} when $\bm{\gamma}_0^*=\bm{Q}\bm{\gamma}_0$ and $\bm{\xi}_0^*=\bm{Q}\bm{\xi}_0$. Individual factor coordinates transform with the basis. Even for fixed $K$, the magnitude of $\bm{\gamma}_0=\bm{B}^\top\bm{\beta}_0$ in $\mathcal M_0$ can grow with the sparse signal.

The error is independent of $\bm{D}_i$ and has a common distribution function $F_\varepsilon$, defining a location-shift model with unrestricted location and no mean or variance requirement. If $F_\varepsilon^{-1}(\tau)$ denotes its generalized quantile, the model gives
\begin{equation}
 Q_{T\mid \bm{D}_i}(\tau)
 =\exp\{\eta_i^0+F_\varepsilon^{-1}(\tau)\},\qquad 0<\tau<1.
 \label{eq:quantile-interpretation}
\end{equation}
Within covariate support, a unit change in $u_{ij}$ at fixed remaining coordinates multiplies conditional event-time quantiles by $\exp(\beta_{0j})$. The coefficients describe conditional log-time associations: larger $\eta_i^0$ corresponds to longer event times. The intercept is absorbed into $F_\varepsilon$ and cancels from pairwise residual comparisons.

\subsection{Censoring and observable local event information}
\label{subsec:censoring}
\begin{assumption}[Common error and conditional censoring]
\label{ass:common-error}
The triples $(\bm{D}_i,\varepsilon_i,C_i)$ are independent and identically distributed, and
\begin{equation}
 \varepsilon_i\perp \bm{D}_i,\qquad
 \varepsilon_i\perp C_i\mid \bm{D}_i.
 \label{eq:censoring-independence}
\end{equation}
The common error distribution is absolutely continuous, with density $f_\varepsilon$ satisfying $f_\varepsilon(t)\le L_\varepsilon<\infty$ for almost every $t\in\mathbb R$.
\end{assumption}
Assumption~\ref{ass:common-error} yields $\Pr(\varepsilon_i\le t\mid \bm{D}_i,C_i)=F_\varepsilon(t)$ and $T_i\perp C_i\mid \bm{D}_i$. Censoring may depend on any covariate block and may include atoms, such as an administrative end of follow-up. The density condition applies to the event error. Recorded-time rounding is handled by a separately specified observation or tie convention.

In $\mathcal M_0$, $\eta_i^0$ is measurable with respect to $(\bm{X}_i,\bm{Z}_i)$, and Assumption~\ref{ass:common-error} also gives $T_i\perp C_i\mid(\bm{X}_i,\bm{Z}_i)$ because $\varepsilon_i$ is independent of $(\bm{D}_i,C_i)$. For $\mathcal M_1$, conditioning on the full design retains the association that can arise between the additional term $\bm{f}_i^\top\bm{\xi}_0$ and censoring after the factors are integrated out.

Define the residual-scale quantities
\begin{equation}
 \begin{gathered}
C_i^*=\log C_i-\eta_i^0,\\
 R_i=Y_i-\eta_i^0=\min(\varepsilon_i,C_i^*),\\
 G(t\mid \bm{D}_i)=\Pr(C_i^*\ge t\mid \bm{D}_i).
 \end{gathered}
 \label{eq:residual-censoring}
\end{equation}
Let $S_\varepsilon(t)=\Pr(\varepsilon_i\ge t)$.

\begin{assumption}[A common local information interval]
\label{ass:overlap}
There exist $t_\star\in\mathbb R$, $A>0$, $m_\varepsilon>0$, and $g_0>0$ such that, on $I=[t_\star-A,t_\star+A]$,
\begin{equation}
 f_\varepsilon(t)\ge m_\varepsilon\quad\text{a.e.},\qquad
 \inf_{t\in I}G(t\mid \bm{D}_i)\ge g_0\quad\text{a.s.}
 \label{eq:local-overlap}
\end{equation}
The constants used for uniform bounds are independent of $n$ and $p$.
\end{assumption}
Assumption~\ref{ass:overlap} supplies a common residual-time interval with positive event density and follow-up probability uniformly over the design, while allowing arbitrary error location. This design-uniform overlap is stronger than control of the marginal censoring proportion. For example, fixed administrative follow-up can violate the overlap condition when $\eta_i^0$ is unbounded above.

Conditioning on $\mathcal \bm{D}_n=\sigma(\bm{D}_1,\ldots,\bm{D}_n)$ fixes factor estimates obtained from the covariates. Patient outcomes remain conditionally independent with potentially different censoring laws. Section~\ref{sec:theory} uses this conditional structure and treats patients as the independent sampling units.

\subsection{Identification of the slope target}
\label{subsec:identification}
Slope uniqueness is governed by the following sparse-design separation condition. It complements the coordinate normalization in Section~\ref{subsec:factors}.

\begin{assumption}[No sparse affine alias]
\label{ass:separation}
For some $\kappa_{\mathrm{pop}}>0$,
\begin{equation}
 \operatorname{Var}(\bm{D}_i^\top \bm{v})\ge\kappa_{\mathrm{pop}}\|\bm{v}\|_2^2
 \quad\text{whenever}\quad
 \bm{v}=(\bm{v}_a^\top,\bm{v}_\beta^\top)^\top,\quad \|\bm{v}_\beta\|_0\le 2s.
 \label{eq:population-separation}
\end{equation}
\end{assumption}
The unrestricted block $\bm{v}_a$ allows joint variation in clinical and factor effects. Assumption~\ref{ass:separation} separates a sparse molecular signal from these effects; a positive lower eigenvalue of $\operatorname{Cov}(\bm{D}_i)$ is sufficient. For $\mathcal M_0$, the analogous restriction on $(\bm{Z}_i,\bm{X}_i)$ suffices. This condition identifies the population target, while Section~\ref{sec:theory} introduces the estimated-design cone condition for penalized estimation.

To state the target precisely without assuming integrable log-time errors, consider two independent patients, indexed by 1 and 2. Put
\[
 e_i^o(\bm{\vartheta})=Y_i-\bm{D}_i^\top\bm{\vartheta},\qquad
 \Theta_s=\{\bm{\vartheta}=(\bm{a}^\top,\bm{\beta}^\top)^\top:\|\bm{\beta}\|_0\le s\}.
\]
The superscript $o$ denotes use of the latent design. Define the expected Gehan \emph{contrast}
\begin{equation}
 \mathcal G(\bm{\vartheta},\bm{\vartheta}_0)
 =E\!\left[\Delta_1
 \left\{\bigl(e_2^o(\bm{\vartheta})-e_1^o(\bm{\vartheta})\bigr)_+
        -\bigl(e_2^o(\bm{\vartheta}_0)-e_1^o(\bm{\vartheta}_0)\bigr)_+\right\}\right].
 \label{eq:oracle-contrast}
\end{equation}
The subtraction is inside the expectation. The 1-Lipschitz positive-part function bounds the integrand by $|(\bm{D}_1-\bm{D}_2)^\top(\bm{\vartheta}-\bm{\vartheta}_0)|$, which is integrable under the covariate moment condition. Thus the Gehan contrast is well defined without an integrable log-time error \citep{CaiEtAl2009,SuderMolstad2022}.

\begin{proposition}[Separation of the rank target]
\label{prop:identification}
Under Assumptions~\ref{ass:common-error}--\ref{ass:separation}, let
\[
 c_{\mathrm{loc}}=A m_\varepsilon^2g_0^2,\qquad
 \rho_A(t)=
 \begin{cases}
 t^2,&|t|\le A,\\
 2A|t|-A^2,&|t|>A.
 \end{cases}
\]
For every $\bm{\vartheta}\in\Theta_s$, with $\bm{v}=\bm{\vartheta}-\bm{\vartheta}_0$,
\begin{equation}
 \mathcal G(\bm{\vartheta},\bm{\vartheta}_0)
 \ge \frac{c_{\mathrm{loc}}}{2}
 E\!\left[\rho_A\{(\bm{D}_1-\bm{D}_2)^\top \bm{v}\}\right].
 \label{eq:population-gap}
\end{equation}
Consequently $\bm{\vartheta}_0$ is the unique minimizer of this contrast over $\Theta_s$ for the fixed latent representation. In $\mathcal M_0$, the analogous contrast formed with residuals $Y_i-\bm{Z}_i^\top\bm{\theta}-\bm{X}_i^\top\bm{\beta}$ identifies $(\bm{\theta}_0,\bm{\beta}_0)$ from the observed-data law, under the corresponding observed-design separation condition.
\end{proposition}

A proof is given in Appendix~\ref{app:identification}. The lower bound uses $\rho_A$, which is quadratic for $|t|\le A$ and linear beyond that range. Section~\ref{sec:theory} develops the corresponding finite-sample bounds.

The primary observed-data target is $(\bm{\theta}_0,\bm{\beta}_0)$ in $\mathcal M_0$. Extension results for $\mathcal M_1$ apply to a specified factor-model sequence with explicit recovery conditions; identification across unrestricted latent decompositions requires further restrictions.

\subsection{Centering and error criteria}
\label{subsec:centering}
Let $\bm{X},\bm{F},\bm{U},\bm{Z}$ collect row vectors $\bm{X}_i^\top,\bm{f}_i^\top,\bm{u}_i^\top,\bm{Z}_i^\top$, respectively, and set
\begin{equation}
 \begin{gathered}
\bm{J}=\bm{I}_n-\mathbf 1_n\mathbf 1_n^\top/n,\\
 \bm{X}_c=\bm{J}\bm{X},\quad \bm{F}_c=\bm{J}\bm{F},\quad \bm{U}_c=\bm{J}\bm{U},\quad \bm{Z}_c=\bm{J}\bm{Z}.
 \end{gathered}
 \label{eq:sample-centering}
\end{equation}
Then $\bm{X}_c=\bm{F}_c\bm{B}^\top+\bm{U}_c$, and
\begin{equation}
 \begin{gathered}
\bm{\eta}_c^0:=\bm{J}\bm{\eta}^0
   =\bm{Z}_c\bm{\theta}_0+\bm{U}_c\bm{\beta}_0+\bm{F}_c\bm{\gamma}_0,
 \\
 Y_i-(\eta_i^0-\bar\eta^0)=R_i+\bar\eta^0.
 \end{gathered}
 \label{eq:centered-predictor}
\end{equation}
The common constant $\bar\eta^0$ cancels from each pairwise residual difference. Sample covariate centering therefore preserves the rank comparisons and slope target while accommodating arbitrary error location.

The primary error quantities are $\|\widehat{\bm{\beta}}-\bm{\beta}_0\|_2$ and $\|\widehat{\bm{\beta}}-\bm{\beta}_0\|_1$. The separate quantity $n^{-1/2}\|J(\widehat{\bm{\eta}}-\bm{\eta}^0)\|_2$ measures location-invariant prediction error on the training sample. Section~\ref{sec:estimation} defines the feasible criterion and reference coordinates for these targets.
